\crefalias{section}{appendix}
\appendix

\section{评估套件规格}
\label{app:suite}

\subsection{参考与度量}
\begin{itemize}
\item 参考：Mixbox~2.0 \texttt{lerpFloat(base, tint, a)}，从官方分发加载；仅供研究对照，不属于算法本身。
\item 只使用浮点路径（$[0,1]$ 的 sRGB 显示值），避免 8 位量化。
\item 度量：CIEDE2000，取 $k_{L}=k_{C}=k_{H}=1$，包含 $G$ 因子、色相角回绕、权重函数 $S_{L},S_{C},S_{H}$ 与旋转项 $R_{T}$。
\item 复刻校验：实现对照 Sharma 等的 34 组参考对验证，最大偏差 $4.95\times10^{-5}$。
\item 置信区间：bootstrap 95\% 百分位区间（1000 次重采样），重采样流固定。
\end{itemize}

\subsection{测试集}
\begin{itemize}
\item T1：调色板全部对（10 个基色 $\times$ 8 个染色）在 $\alpha\in\{0.25,0.5,0.75\}$（240 个样本）。
\item T2：19 个压力对在三个配比下（57 个样本），列于 \cref{tab:t2}。
\item T3：均匀随机 sRGB 对，$\alpha$ 均匀，seed 20260808，$n=2000$（可配置）。
\item T4：10 个颜色在 $\alpha\in\{0,1\}$ 下自混合；通道容差 0.01。
\item T5：20 个固定深色与 24 个随机深色，各在 $\alpha\in\{0.1,0.2,\dots,0.9\}$（396 个样本）。
\item T5-留出：seed 20260812 的 24 个随机深色，同一 $\alpha$ 网格（216 个样本）。
\item tail：20 个 T5 固定色在 $\alpha\in\{0.92,0.93,\dots,0.99\}$（160 个）与 9 个盲区色在 T5 $\alpha$ 网格上（81 个）。
\item T5E：20 个 T5 固定色在深重侧 $\alpha$（120 个）、12 个近中性色在深重侧 $\alpha$（72 个）以及同 12 色在浅重侧 $\alpha$（48 个）。
\end{itemize}

\begin{table}[t]
\centering
\footnotesize
\setlength{\tabcolsep}{3pt}
\caption{T2 压力对。}
\label{tab:t2}
\begin{tabular}{p{0.16\linewidth}p{0.78\linewidth}}
\toprule
类别 & 颜色对 \\
\midrule
互补对（3） & \#E8321F + \#00D5C0；\#FFD61A + \#2E5EFF；\#E33BA8 + \#00B359 \\
RYB 错配（6） & \#E8321F + \#00B359；\#2E5EFF + \#FF7A1A；\#E8321F + \#2E5EFF；\#FFD61A + \#00B359；\#FFD61A + \#00D5C0；\#E33BA8 + \#00D5C0 \\
白黑混合（6） & \#E8321F、\#2E5EFF、\#FFD61A，各 $\times$ \{\#FFFFFF, \#000000\} \\
相似色（4） & \#E8321F + \#FF7A1A；\#FFD61A + \#FF7A1A；\#2E5EFF + \#8A5BE0；\#00B359 + \#00D5C0 \\
\bottomrule
\end{tabular}
\end{table}

\subsection{T5 与 T5E 套件}
\label{app:t5}

\paragraph{T5 固定色.}
\Cref{tab:t5colors} 列出 20 个固定色：10 个痛点色与 10 个家族扩展色。随机深色用拒绝采样直到 $\max(\mathbf{c})\in[0.30,0.60]$、$\min(\mathbf{c})\le0.08$、$\max(\mathbf{c})-\min(\mathbf{c})\ge0.25$；T5 用 seed 20260810，T5-留出用 seed 20260812。

\begin{table}[t]
\centering
\small
\caption{T5 固定深色（sRGB 十六进制与显示值浮点）。}
\label{tab:t5colors}
\begin{tabular}{llll}
\toprule
Hex & 浮点 RGB & Hex & 浮点 RGB \\
\midrule
$\#0D1B44$ & 0.051, 0.106, 0.267 & $\#0A2A6E$ & 0.039, 0.165, 0.431 \\
$\#000080$ & 0.000, 0.000, 0.502 & $\#8E1F0A$ & 0.557, 0.122, 0.039 \\
$\#001B75$ & 0.000, 0.106, 0.459 & $\#005A32$ & 0.000, 0.353, 0.196 \\
$\#001E80$ & 0.000, 0.118, 0.502 & $\#00552B$ & 0.000, 0.333, 0.169 \\
$\#080248$ & 0.031, 0.008, 0.282 & $\#00585E$ & 0.000, 0.345, 0.369 \\
$\#5F1708$ & 0.373, 0.090, 0.031 & $\#2A0A5E$ & 0.165, 0.039, 0.369 \\
$\#5E0F1A$ & 0.369, 0.059, 0.102 & $\#4A0B6E$ & 0.290, 0.043, 0.431 \\
$\#4A3A05$ & 0.290, 0.227, 0.020 & $\#6E0B4A$ & 0.431, 0.043, 0.290 \\
$\#280501$ & 0.157, 0.020, 0.004 & $\#7A3B00$ & 0.478, 0.231, 0.000 \\
$\#250903$ & 0.145, 0.035, 0.012 & $\#6E4A00$ & 0.431, 0.290, 0.000 \\
\bottomrule
\end{tabular}
\end{table}

\paragraph{T5E.}
深重侧 $\alpha$ 取 0.001、0.005、0.01、0.02、0.05、0.1；浅重侧取 0.92、0.95、0.98、0.99。近中性组为 \#000B2E、\#06052E、\#1A0F22、\#0A1A18，外加 seed 20260816 的 8 个随机色（拒绝采样：$\max\in[0.04,0.10]$，色度范围 $\in[0.02,0.06]$）。

\paragraph{tail 盲区色.}
$\#0D1B44$、$\#001B75$、$\#000080$、$\#001E80$、$\#080248$、$\#0A2A6E$、$\#2A0A5E$、$\#060538$、$\#000333$。

\subsection{划分与统计}
\begin{itemize}
\item 拟合抽样：T1、T2、T5、T5E、tail、T3（seed 20260808）前 800 对、6000 个伪随机对（seed 999），以及 12 个锚点互补色对（各 19 个配比）。
\item 留出：同一次抽样的 T3 第 800--1999 对、seed 20260809 与 20260811 的 T3 抽样、T5-留出（seed 20260812）。
\item 对每个模型与集合报告均值、最大值、中位数、95 分位、标准差，以及（随机集）均值的 bootstrap 95\% 区间。
\item 图表由仓库工具从参考实现生成；每个抽样的 seed 在图表说明或本附录中给出。
\end{itemize}

\section{谐波修正系数}
\label{app:coeff}

\cref{eq:gamma} 的 51 个系数按特征顺序列于下（\cref{tab:features}：截距、\texttt{lum}、\texttt{dgug}、\texttt{dcmucm}、\texttt{wUcP}、\texttt{u2}、\texttt{pm}、\texttt{dg}、\texttt{op}、\texttt{span}、\texttt{drur}、\texttt{pn}、\texttt{Mmn}、\texttt{lumU}、\texttt{ur}、\texttt{dcAbs}、\texttt{gRu}）；截距在首位。

\paragraph{R 通道（17 个系数）.}
$-0.24945965$、$-0.00458899$、$0.00415136$、$-0.03758060$、$0.05450570$、$0.09989345$、$0.25226195$、$0.02321094$、$0.15029394$、$0.03230901$、$0.03005819$、$0.34621807$、$-0.57046216$、$-0.00274566$、$0.03208997$、$-0.05313629$、$-0.02941071$。

\paragraph{G 通道（17 个系数）.}
$-0.15650537$、$0.14061150$、$-0.03620353$、$0.03735064$、$0.02174583$、$0.02381613$、$-0.06481304$、$-0.02027288$、$0.13851721$、$0.13327492$、$0.01269829$、$0.03768178$、$-0.05558582$、$0.03768944$、$0.01505410$、$0.02180116$、$-0.01767582$。

\paragraph{B 通道（17 个系数）.}
$0.16757082$、$-0.04373748$、$0.02221282$、$-0.02131140$、$-0.08893948$、$0.05973609$、$-0.09324463$、$-0.00242072$、$-0.11889782$、$-0.08066689$、$-0.01881345$、$-0.40457567$、$0.51586855$、$-0.01855359$、$0.03390971$、$-0.01807147$、$-0.04148750$。

\section{合成光谱数据集}
\label{app:synthetic}

\subsection{数据集}
模型为两常数 K--M，每个颜料在 380--750 nm 的 10 nm 网格（38 个采样）上分别给出吸收与散射：
\begin{equation}
K(\lambda)=\kappa_{\mathrm{base}}+\sum_{i}A_{i}e^{-((\lambda-\mu_{i})/\sigma_{i})^{2}}+\sum_{j}\frac{B_{j}}{1+e^{\pm((\lambda-\lambda_{j})/w_{j})}},
\qquad S(\lambda)=S_{\mathrm{base}}.
\end{equation}
每个参数都是固定常数；数据集由工具确定性复现。带位置依据第一性原理颜料化学：黄、橙、红与土族在蓝—绿区带陡峭 logistic 吸收边（$\lambda_{0}\approx490$--$585$ nm），蓝与绿在红侧带上升边（$\approx610$--$660$ nm），品红与紫在绿区带高斯带（$\approx550$--$580$ nm），过渡金属或有机颜料另加特征带（$\approx430$--$660$ nm）。幅度、宽度、基底与散射由坐标下降标定到 BabelColor X-Rite ColorChecker Classic Avg30 光谱 \cite{pascale2012}（拟合 $\Delta E_{00}\le1.2$，多数色卡 $\le0.2$）。这些色卡是配方色卡光谱而非纯颜料厚涂光谱；K--M 方程对任意反射率光谱都是封闭的，且该色卡面向 RGB 显示设计，故无需显示标定。

数据集包含 29 个颜料：24 个功能色相（白、黑、米色、浅/中/深灰、蓝灰、两黄、两橙、两红、两品红、紫、三蓝、青、两绿、两土色）与 5 个深色派生（深蓝、红、绿、品红、紫），用于 P5 套件。仓库不含任何第三方测量数据；数值参数离线拟合自公开来源，拟合值随文发表。

\subsection{前向管线与验证}
前向管线分别叠加 $K=(1-\alpha)K_{A}+\alpha K_{B}$ 与 $S=(1-\alpha)S_{A}+\alpha S_{B}$，求 $\kappa=K/S$ 与 \cref{eq:km}，在 10 nm 网格上对 D65 与 CIE 1931 $2^{\circ}$ 观察者积分（梯形法，归一化到 $Y=100$），再转为 IEC 61966-2-1 sRGB，并做截断与色域外标记。

数据集由工具自检：全部 29 个颜料满足 $\Rinf\in(0,1)$ 且 $K\ge0$、$S>0$；纯原色的 unmix 重建 $\Delta E_{00}=0.000$；光谱混合在 $\alpha=0,1$ 处恒等；黄 $\times$ 蓝得到绿（$h^{\circ}\in(110,190)$）；品红 $\times$ 绿与红 $\times$ 青得到浑浊色（$C^{\ast}<30$ 与 $<40$）；深 $+$ 白保持色相（$\alpha=0.25$--$0.75$ 上 $|\Delta h|<40^{\circ}$）；每个颜料都落在其族色相窗内。饱和纯颜料可能合法地超出 sRGB 色域；这些样本被标记并在输入模型前截断（P1：6/828；P3：35/2000；P5：0/135）。

\subsection{套件}
\begin{itemize}
\item P1：24 色相调色板的全部 $\binom{24}{2}$ 个对，$\alpha\in\{0.25,0.5,0.75\}$（828 个样本）。
\item P2：九个互补/压力对，同一组配比（27 个样本）。
\item P3：2000 个随机颜料对，$\alpha$ 均匀，seed 20260813。
\item P5：15 个深色颜料 $\times$ 白，$\alpha\in\{0.1,\dots,0.9\}$（135 个样本）。
\item P5-buff：同一深色集 $\times$ 米色，同一网格（135 个样本）。
\item P5E：深色集 $\times$ 白，$\alpha\in\{0.001,0.005,0.01,0.02,0.05,0.1,0.92,0.95,0.98,0.99\}$（150 个样本）。
\end{itemize}

\section{反解测量}
\label{app:inversion}

本附录中的全部数字都在 Mixbox~2.0 参考（浮点路径）上测得。

\subsection{反解幂指数}
对锚点色，按 \cref{eq:inversion} 逐通道反解指数。\Cref{tab:d1} 报告合成锚点色吸收带在 $\alpha\in\{0.1,\dots,0.9\}$ 上的反解指数（反解最大通道）。各族漂移均小于 0.2；$\alpha\to1$ 处的崩坏正是 $\varphi$ 映射处理的范围。

\begin{table}[t]
\centering
\small
\caption{吸收带的反解指数 $p$ 随 $\alpha$ 的变化（实测）。}
\label{tab:d1}
\begin{tabular}{lccccc}
\toprule
族 & $\alpha=0.1$ & $\alpha=0.3$ & $\alpha=0.5$ & $\alpha=0.7$ & $\alpha=0.9$ \\
\midrule
$B$ & 0.97 & 0.97 & 0.98 & 1.00 & 1.03 \\
$R$ & 0.98 & 1.00 & 1.04 & 1.09 & 1.16 \\
$G$ & 1.00 & 1.03 & 1.07 & 1.12 & 1.18 \\
\bottomrule
\end{tabular}
\end{table}

\Cref{tab:d2} 是 $\alpha=0.5$ 处锚点反解表的摘录，其中 $s$ 为通道沿色度方向的位置。对合成色 $(0.01,v,0.01)$，$s=1$ 通道（G 最大）给出 $v=0.05\rightarrow1.07$、$0.12\rightarrow1.06$、$0.25\rightarrow0.81$、$0.35\rightarrow0.58$、$0.55\rightarrow0.42$；快速亮端区仅在 $v>0.25$ 时适用。

\begin{table}[t]
\centering
\small
\caption{锚点反解表（摘录，$\alpha=0.5$）。}
\label{tab:d2}
\begin{tabular}{lccccc}
\toprule
颜色 & 通道 & $v$ & $s$ & $p$ \\
\midrule
$\#0D1B44$ & R & 0.051 & 0.00 & 2.05 \\
$\#0D1B44$ & G & 0.106 & 0.26 & 0.84 \\
$\#0D1B44$ & B & 0.267 & 1.00 & 0.23 \\
$\#080248$ & R & 0.031 & 0.09 & 2.25 \\
$\#080248$ & B & 0.282 & 1.00 & 0.12 \\
$\#4A0B6E$ & B & 0.431 & 1.00 & 0.85 \\
$\#6E0B4A$ & G & 0.043 & 0.00 & 1.59 \\
$\#5F1708$ & R & 0.373 & 1.00 & 0.76 \\
$\#5F1708$ & G & 0.090 & 0.17 & 1.77 \\
$\#5F1708$ & B & 0.031 & 0.00 & 0.90 \\
$\#8E1F0A$ & R & 0.557 & 1.00 & 0.62 \\
$\#005A32$ & G & 0.353 & 1.00 & 0.57 \\
$\#005A32$ & B & 0.196 & 0.56 & 1.04 \\
$\#6E0B4A$ & B & 0.290 & 0.64 & 0.91 \\
\bottomrule
\end{tabular}
\end{table}

\subsection{残差与耦合}
对 $v=0$ 的通道，参考曲线与颜料无关、只依赖通道（例如 $\#000080$ 的 R 与 G 都为零，但参考曲线不同），这引出残差项。拟合中测得的残差约为 1.85（R）、1.0（G）、1.1（B），并随通道值衰减微调。

耦合项通过固定 $B=0.431$、$G=0.043$ 并在 $\alpha=0.5$ 处扫 $R$ 测得（\cref{tab:d3}）。关系是单调的 S 形：紫 $\#4A0B6E$（$R=0.29$）的 $p_{B}=0.85$，而纯蓝（$R\approx0$）的 $p_{B}=0.12$--$0.17$。耦合是对称的：高 $B$ 使 $R$ 变慢（$B$ 从 0 增至 0.431 时 $p_{R}$ 从 1.04 升到 2.02）。

\begin{table}[t]
\centering
\small
\caption{通道耦合：反解 $p_{B}$ 随 $R$ 的变化（实测）。}
\label{tab:d3}
\begin{tabular}{lccccccc}
\toprule
$R$ & 0.03 & 0.06 & 0.10 & 0.15 & 0.20 & 0.25 & 0.30 \\
$p_{B}$ & 0.17 & 0.21 & 0.35 & 0.60 & 0.73 & 0.80 & 0.85 \\
\bottomrule
\end{tabular}
\end{table}

\subsection{方向响应}
逐 $\alpha$ 反解 T5 红主导固定色可保持 $\alpha$ 恒定性（$\sigma<0.1$）；\cref{tab:d4} 报告逐通道需求。出现三个结构信号：红族 G 保持（$\#5E0F1A$ 与 $\#8E1F0A$ 需要 $G\approx1.7$--$1.8$）、黄褐 B 保持（$\#906D00$ 需要 $B\approx1.5$）、红透射（R 需求为 0.62--0.77）。对 11 个红主导固定色的反解残差做最小二乘（三通道同一线性形式 $\delta_{c}=a+k\,v$），得到 \cref{eq:deltas} 的统一响应。范围门 $g_{R}\sigma_{8}(v_{R}-\mathrm{RTH})$（$\mathrm{RTH}=0.21929798$）是必要的，因为方向系数在其下改变符号。

\begin{table}[t]
\centering
\small
\caption{红族反解：逐通道指数需求（实测）。}
\label{tab:d4}
\begin{tabular}{lccc}
\toprule
颜色 & R 需求 & G 需求 & B 需求 \\
\midrule
$\#5E0F1A$ 深红 & 0.77 & 1.78 & 0.89 \\
$\#8E1F0A$ 深红 & 0.62 & 1.71 & 0.88 \\
$\#906D00$ 赭黄 & 0.72 & 0.95 & 1.49 \\
$\#940C7C$ 品红 & 0.99 & 1.31 & 0.98 \\
\bottomrule
\end{tabular}
\end{table}

\subsection{互补族}
在互补色对上反解参考可见族结构化的残差：对每个色相对立族，误差集中在混合区间中部，形状为 $4a(1-a)$，幅度随两个输入颜色平滑变化。例如亮黄 $\times$ 深蓝族的 G 通道误差近零、B 通道误差至多 0.05，而 R 通道误差在 $\alpha\in[0.4,0.9]$ 上按 $4a(1-a)$ 变化、幅度 0.28--0.29。由于参考的潜空间非线性，幅度跨族平滑变化，单一幅值不足；因此谐波修正是十七特征模板上的线性模型，在互补族语料上拟合（权重 10），每通道 17 个系数。端点抑制 $h(a)$ 与色度衰减 $\mathrm{cSat}$ 是结构性而非拟合得到的。
